\section{Experiment}
%这个是作为附加组件的对比

\begin{table*}[t]
\centering
\small
% \setlength{\tabcolsep}{5pt}
% \renewcommand{\arraystretch}{1.15}
% \usepackage{amssymb} % 这里应该说是通用框架
\caption{\small {Comparison with long-video RAG baselines. ``Native Input'' indicates whether the model processes video directly without relying on external auxiliary text or tools. Bold and underlined numbers indicate the best and second-best accuracies, respectively. ‘–’ indicates not applicable.}}
\resizebox{\linewidth}{!}{
\begin{tabular}{lc cc cc cc}
\toprule
\multirow{2}{*}{\textbf{Model}} & \multirow{2}{*}{\textbf{Native Input}} &
\multicolumn{2}{c}{\textbf{LV-Bench (Val)}} &
\multicolumn{2}{c}{\textbf{MLVU}} &
\multicolumn{2}{c}{\textbf{Video-MME-Long (w/o Sub)}} \\
\cmidrule(lr){3-4} \cmidrule(lr){5-6} \cmidrule(lr){7-8}
 & & Overall $\uparrow$ & Gain (\%) $\uparrow$ &
     Overall $\uparrow$ & Gain (\%) $\uparrow$ &
     Overall $\uparrow$ & Gain (\%) $\uparrow$ \\
\midrule
% GPT-4o & \cmark & 66.7 & - & 64.6 & - & 65.3 & - \\
% + Video-RAG& \xmark & 67.2 & 0.7 & 65.6 &1.5 & & \\
% + TV-RAG & \xmark & 67.0 & 0.4&  65.5 & 1.4&&  \\
% \rowcolor[rgb]{0.80,0.90,0.95} \textbf{+ VideoStir (Ours)} & \cmark & \textbf{68.5} & \textbf{2.7} & \textbf{67.0} & \textbf{3.7} &  &  \\ %69.8
% \tightmidrule
LLaVA-Video (7B) & \cmark & 56.6 & - & 70.8 & - & - & - \\
+ Video-RAG & \xmark & 58.7 & 3.7 & \underline{72.4} & \underline{2.3} & - & - \\
+ TV-RAG & \xmark & 58.8 & 2.1& 72.1 & 1.8&- & - \\
\rowcolor[rgb]{0.80,0.90,0.95} \textbf{+ VideoStir (Ours)} & \cmark & \textbf{60.3} & \textbf{6.5} & \textbf{73.1}& \textbf{3.2} & - & - \\ %64.7
\tightmidrule
mPLUG-Owl3 (8B) & \cmark & 52.1 & - & 63.7 & - &50.1 & - \\
+ Video-RAG & \xmark & 54.5   & 4.6 & \underline{64.4} & \underline{1.1}& 50.6 & 1.0 \\
+ TV-RAG & \xmark & 54.8 & 5.2& 64.0 &0.5 & \underline{50.9}& \underline{1.6} \\
\rowcolor[rgb]{0.80,0.90,0.95} \textbf{+ VideoStir (Ours)} & \cmark & \textbf{55.1} & \textbf{5.8} & \textbf{64.6} & \textbf{1.4} & \textbf{51.0} &\textbf{1.8}  \\
\tightmidrule
Aria (25B) & \cmark & 64.2 & - & 70.6 & - & 58.8 & - \\
+ Video-RAG & \xmark & \textbf{66.4} & \textbf{3.4} & \textbf{72.1} & \textbf{2.1} & \textbf{59.6}& \textbf{1.4} \\
+ TV-RAG & \xmark & 65.5 & 2.0& 71.2& 0.8&59.0& 0.3\\
\rowcolor[rgb]{0.80,0.90,0.95} \textbf{+ VideoStir (Ours)} & \cmark & \underline{65.8} & \underline{2.5} &  \underline{71.7} & \underline{1.6} & \underline{59.4} & \underline{1.0} \\
\tightmidrule
InternVL-1.5 (26B) & \cmark & 51.2 & - & 50.4 & - & 45.6 & - \\
+ Video-RAG& \xmark & 52.2 & 2.0 &\textbf{51.8} & \textbf{2.8} & \underline{46.7} & \underline{2.4} \\
+ TV-RAG & \xmark & \underline{52.4} & \underline{2.3}& 51.0 &1.2 & 46.3& 1.5\\
\rowcolor[rgb]{0.80,0.90,0.95} \textbf{+ VideoStir (Ours)} & \cmark & \textbf{52.7} & \textbf{2.9} & \underline{51.6}&\underline{2.4} &\textbf{46.9} & \textbf{2.9} \\ 
\tightmidrule
LLaVA-Video (72B) & \cmark & 61.9 & - & 73.1 & - &61.5 & - \\
+ Video-RAG & \xmark & \underline{65.4} & \underline{5.7} & \underline{73.8} & \underline{1.0} & \textbf{62.3} & \textbf{1.3} \\
+ TV-RAG & \xmark & 64.6 & 4.3 &  73.4& 0.4& 61.8 & 0.5 \\
\rowcolor[rgb]{0.80,0.90,0.95} \textbf{+ VideoStir (Ours)} & \cmark & \textbf{66.0} &\textbf{6.6}&  \textbf{74.1}  & \textbf{1.4} & \underline{62.1} & \underline{1.0} \\ %72.8
\tightbottomrule
\end{tabular}
}
\label{tab:comp1}
% \vspace{-0.10 in}
\end{table*}




\begin{table}[t]
\centering
\small
% \setlength{\tabcolsep}{1pt}
\caption{\small {System-level comparison on EgoSchema's test set with SOTA methods built upon fixed (M)LLMs. These approaches can be categorized into two groups: one employs video-to-text tools (e.g., captioners or object detectors) to provide auxiliary textual inputs for downstream (M)LLMs, while the other leverages the native multimodal understanding capability (Native Input) of MLLMs without auxiliary text. Bold and underlined numbers indicate the best and second-best accuracies, respectively. ‘–’ indicates not applicable.}}
\resizebox{\linewidth}{!}{
\begin{tabular}{l l c}
\toprule
(M)LLM Backbone & Method & Overall$\uparrow$ \\
% \midrule
\specialrule{\lightrulewidth}{0pt}{0pt}
\rowcolor[HTML]{FFF7F0}
\multicolumn{3}{c}{\rule[-0.9ex]{0pt}{3.3ex}\textbf{\textit{Methods with Video-to-Text Tools}}}\\
\specialrule{\lightrulewidth}{0pt}{0.8ex}
\multirow{2}{*}{GPT-3.5} & LLoVi~\cite{llovi} & 57.6 \\
                       & DrVideo~\cite{drvideo}        & 62.6 \\
\midrule
\multirow{4}{*}{GPT-4} & VideoAgent~\cite{videoagent} & 60.2 \\
                       & LLoVi~\cite{llovi}          & 61.2 \\
                       & VidAgent~\cite{vidagent}  & 62.8 \\
                       & VideoTree~\cite{videotree}  & 66.2\\
                       & DrVideo~\cite{drvideo}               & 66.4 \\
\midrule
Qwen2-VL-7B     & Vgent~\cite{vgent}    &  \textbf{68.0}\\
% \midrule
\specialrule{\lightrulewidth}{0pt}{0pt}
\rowcolor[HTML]{FFF7F0}
\multicolumn{3}{c}{\rule[-0.9ex]{0pt}{3.3ex}\textbf{\textit{Methods with Native Input}}}\\
\specialrule{\lightrulewidth}{0pt}{0.8ex}
\multirow{2}{*}{Qwen2-VL-7B} & IG-VLM~\cite{igvlm} & 66.2 \\
                            & VideoStir (Ours)         & \underline{67.2}\\
% \midrule
% \multirow{2}{*}{GPT-4o}      & IG-VLM~\cite{igvlm} & 73.0 \\
%                             &  \textbf{VideoStir (Ours)}         & \textbf{75.8} \\
\bottomrule
\end{tabular}
}
\label{tab:comp2}
% \vspace{-0.10 in}
\end{table}

%指名Table 1的Setting与Video-RAG一致

%Table3的Setting与FOCUS与AKS一致
\subsection{Dataset and Implementation}
\noindent\textbf{Dataset.} We curate IR-600K, a  dataset designed to enhance small-scale MLLMs’ ability to judge intent relevance between video frames and queries, supporting intent-aware long-video RAG. To the best of our knowledge, IR-600K is the first dataset targeting intent-level frame-query alignment, i.e., whether a frame provides evidence that satisfies the query’s underlying information need. This differs from prior datasets such as TVSum~\cite{tvs} and QVHighlights~\cite{qah}, which primarily evaluate the semantic saliency of frames with respect to descriptive statements. 

Specifically, we sample 4,678 videos from three representative video QA datasets (ActivityNet-QA~\cite{A-QA}, NExT-QA~\cite{N-QA}, and STAR~\cite{S-QA}).
%to cover diverse tasks including scene recognition, causal/temporal/logical reasoning, and visual description. 
For each video, we uniformly extract frames at 3~fps to form query--frame pairs. We then distill supervision from a powerful MLLM: Qwen2.5-VL-72B-Instruct~\cite{bai2025qwen2} serves as the teacher to produce frame-level intent-relevance labels. Guided by a carefully designed system prompt specifying explicit intent-relevance criteria, the teacher assesses intent-level alignment between each frame and the query, yielding intent-relevance scores. IR-600K comprises 605,676 training samples and 21,791 validation samples. Further details on the intent-relevance scoring prompt, dataset statistics and annotation format are provided in Appendix~\ref{appendix:dataset} and~\ref{appendix:prompt}. 


\noindent\textbf{Implementation.} We adopt Qwen2.5-VL-3B-Instruct~\cite{bai2025qwen2} as the backbone of the intent-relevance scorer and fine-tune it with LoRA adapters with a rank of 16, a scaling factor of 32, and a dropout of 0.05. We optimize the LoRA using AdamW with a learning rate of $5\times10^{-5}$ and a weight decay of 0.05, together with a cosine learning-rate schedule and a warm-up ratio of 0.05. Training is performed for 1 epoch with a batch size of 128 under bf16 precision. For the event boundary detector, we use the vision tower from the Qwen2.5-VL series as the embedding backbone. For spatio-temporal graph construction and multi-hop retrieval, we use the Perception Encoder~\cite{perceptionenc} as the video-language encoder $E_{\mathrm{vl}}$ and set $N=3$, $L=2$, $\eta=0.4$, and $\kappa_s=3.25$; an ablation study on these retrieval hyperparameters is provided in Appendix~\ref{appendix:hyperparam}. All training runs  and comparative experiments are conducted on 8$\times$A100 GPUs, while the remaining experiments are conducted on a single A100 GPU. Unless otherwise specified, results are reported as the median of three trials, and ablations are conducted on LLaVA-Video (7B) using LongVideoBench~\cite{lvb}.

\subsection{Comparison Against Other Methods} %引用一下对比表格里面出现的MLLMs
Table~\ref{tab:comp1} compares VideoStir with the native MLLM and state-of-the-art (SOTA) long-video RAG baselines on LongVideoBench (LV-Bench)~\cite{lvb}, MLVU~\cite{mlvu}, and Video-MME-Long~\cite{videomme}, following the same experimental setup as Video-RAG~\cite{Video-RAG} and TV-RAG~\cite{TVRAG}. Baseline MLLMs include GPT-4o~\cite{hurst2024gpt}, LLaVA-Video (7B/72B)~\cite{zhang2024video}, mPLUG-Owl3 (8B)~\cite{ye2024mplug}, Aria (25B)~\cite{aria}, and InternVL-1.5 (26B)~\cite{chen2024far}. Results are compiled from the benchmarks’ leaderboards and the original papers, with the remaining results obtained from our reproductions based on their open-sourced codes and original settings. In most cases, VideoStir outperforms both the native model and prior long-video RAG baselines, demonstrating highly competitive performance. Notably, VideoStir preserves the native MLLM input setting by operating solely on retrieved visual evidence: it introduces no auxiliary text, tool calls, or extra modalities, while still achieving SOTA performance.  

Table~\ref{tab:comp2} compares VideoStir on EgoSchema against specialized long-video systems built atop fixed (M)LLMs~\cite{egoschema}; baseline results are taken from DrVideo~\cite{drvideo} and Vgent~\cite{vgent}. Relying solely on native multimodal inputs without external video-to-text tools (e.g., captioners or object detectors), VideoStir achieves performance comparable to these specialized systems that leverage such auxiliary tools.

% Table~\ref{tab:comp3} compares VideoStir with long-video keyframe selection baselines on LongVideoBench and Video-MME-Long under the same experimental setting as prior work; baseline results are taken from FOCUS~\cite{FOCUS}. VideoStir achieves competitive performance across evaluated benchmarks. These results demonstrate VideoStir’s effectiveness in improving downstream MLLMs for long-video understanding.

% % ===== Table 1: 组件级消融 =====
% \begin{table}[t]
% \centering
% \small
% \setlength{\tabcolsep}{6pt}
% \renewcommand{\arraystretch}{1.12}
% \caption{Component ablations. Metrics reported on \textit{[Benchmark]}: Overall(\%), \#Text per case, and end-to-end Latency.}
% \begin{tabular}{lccc}
% \toprule
% \textbf{Method Variant} & \textbf{Overall$\uparrow$} & \textbf{\#Text$\downarrow$} & \textbf{Latency(ms)$\downarrow$}\\
% \midrule
% \multicolumn{4}{l}{\emph{w/o Reranker: similarity-based retrieval only}}\\
% \quad CLIP (image)                     & -- & -- & -- \\
% \quad Video-CLIP (video)               & -- & -- & -- \\
% \quad PE (project-specific encoder)    & -- & -- & -- \\
% \midrule
% w/o weighted expectation (use argmax label) & -- & -- & -- \\
% w/o semantic-adaptive segmentation (fixed window graph) & -- & -- & -- \\
% w/o spatial edges                       & -- & -- & -- \\
% w/o temporal edges                      & -- & -- & -- \\
% w/o spatio-temporal graph (uniform sparse frames $\rightarrow R_{\theta}$) & -- & -- & -- \\
% \textbf{Full (STC-R$^{2}$)}             & \textbf{--} & \textbf{--} & \textbf{--} \\
% \bottomrule
% \end{tabular}
% \end{table}


% \begin{table}[t]
% \centering
% \scriptsize
% \setlength{\tabcolsep}{2pt}
% \caption{\small Comparison with keyframe selection methods for long-video understanding. Results are obtained with LLaVA-Video (7B). Bold and underlined numbers indicate the best and second-best accuracies, respectively. ‘–’ indicates not applicable.} 
% \resizebox{\linewidth}{!}{
% \begin{tabular}{l c c c c}
% \toprule
% \multirow{2}{*}{\textbf{Method}}&\multicolumn{2}{c}{\textbf{LV-Bench (Val)}}&\multicolumn{2}{c}{\textbf{Video-MME-Long (w/o Sub)}}\\ \cmidrule(lr){2-3}  \cmidrule(lr){4-5} 
%  & Overall$\uparrow$ & Gain$\uparrow$&Overall$\uparrow$ & Gain$\uparrow$\\
% \midrule
% Uniform Sampling &56.6 &- &54.3&-\\
% AKS~\cite{AKS}&62.1 & 9.7&54.7&0.7\\
% FOCUS~\cite{FOCUS}& 63.5&12.2 &\textbf{56.1}&\textbf{3.3}\\
% \textbf{VideoStir (Ours)}& \textbf{64.3} & \textbf{13.6}&\underline{55.9}&\underline{2.9}\\
% \bottomrule
% \end{tabular}
% }
% \label{tab:comp3}
% % \vspace{-0.1 in}
% \end{table}

% ===== Table 1: 组件级消融 =====

\begin{table}[t]
\centering
\small
\setlength{\tabcolsep}{3pt}
\caption{\small Component-level ablations on the expanded 1k validation set. ``Overall'' denotes overall accuracy, and ``Retrieval Acc.'' measures the accuracy of successfully retrieving key cues. ``w/o Intent-relevance Scorer'' replaces the scorer with similarity-based frame retrieval; ``w/o Weighted Expectation'' replaces probability-weighted expectation with discrete scores; ``w/o Spatio-Temporal Graph'' disables structured retrieval, which bypasses the graph and directly selects the top-14 clips by PE-embedding similarity, following the best-performing setting in Table~\ref{tab:ablation_hyperparams}. Bold indicates the best results.}
\resizebox{\linewidth}{!}{
\begin{tabular}{lccc}
\toprule
\textbf{Method} & \textbf{Overall$\uparrow$} & \textbf{Retrieval Acc.$\uparrow$}\\
\midrule
w/o Intent-relevance Scorer \\
\quad + CLIP L/14                   & 52.8 & 69.2 \\
\quad + Video-CLIP                  & 56.3 & 75.4 \\
\quad + PE                          & 58.1 & 79.8 \\
w/o Weighted Expectation            & 54.2 & 71.6 \\
w/o Event-Boundary Detector         & 62.3 & 88.2 \\
w/o Spatio-Temporal Graph           & 56.4 & 74.8 \\
w/o Spatial Edges                   & 57.2 & 79.3 \\
w/o Temporal Edges                  & 59.8 & 83.4 \\
\midrule
\textbf{Full}                 & \textbf{64.5} & \textbf{92.2}  \\
\bottomrule
\end{tabular}
}
\label{tab:ab1}
% \vspace{-0.10 in}
\end{table}

\begin{table}[t]
\centering
\small
\setlength{\tabcolsep}{6pt}
\renewcommand{\arraystretch}{1.12}
\caption{\small 
Ablations of VideoStir's video-language embedding backbones with $N=1000$. ``Overall'' denotes overall accuracy, and ``Retrieval Acc.'' refers to the accuracy of successfully retrieving the key evidence for answering the query. The bold numbers represent the best results.} 
\resizebox{\linewidth}{!}{
\begin{tabular}{lcc}
\toprule
\textbf{Embedding Backbone} & \textbf{Overall$\uparrow$} & \textbf{Retrieval Acc.$\uparrow$} \\
\midrule
Video-CLIP~\cite{videoclip}              & 62.1 & 88.3  \\
X-CLIP~\cite{xclip}                      & 62.8 & 89.4 \\
\textbf{PE}~\cite{perceptionenc}         & \textbf{64.5} & \textbf{92.2}  \\
\bottomrule
\end{tabular}
}
\label{tab:ab2}
% \vspace{-0.10 in}
\end{table}

\begin{table*}[t]
\centering
\small
\caption{\small {Ablations on the intent-relevant scorer’s distillation strategies. ``IR-600K CE'' and ``LV-Bench CE'' denote the cross-entropy loss computed on the corresponding datasets. To reduce distributional bias, the student model is distilled using teacher models from the same family. ``Retrieval Acc.'' refers to the accuracy of successfully retrieving the key evidence.}} 
\resizebox{\linewidth}{!}{
\begin{tabular}{llccccc}
\toprule
\textbf{Base Model} &\textbf{Method} & \textbf{IR-600K CE $\downarrow$}& \textbf{LV-Bench CE $\downarrow$}&\textbf{Overall $\uparrow$}&\textbf{Retrieval Acc. $\uparrow$}& \textbf{Trainable Params $\downarrow$}\\
\specialrule{\lightrulewidth}{0pt}{0pt}
\rowcolor[HTML]{FFF7F0}
\multicolumn{7}{c}{\rule[-1ex]{0pt}{3.3ex}\textbf{\textit{Teacher Models}}}\\
\specialrule{\lightrulewidth}{0pt}{0.8ex}
InternVL2.5-78B&Zero-Shot   &  -&-& 64.8 & 92.6 &  - \\
Qwen2.5-VL-72B&Zero-Shot   & - &-& 65.4 & 95.8 &  - \\
Qwen2-VL-72B&Zero-Shot   & -&-& 64.9 & 94.7 &  - \\
\specialrule{\lightrulewidth}{0pt}{0pt}
\rowcolor[HTML]{FFF7F0}
\multicolumn{7}{c}{\rule[-1ex]{0pt}{3.3ex}\textbf{\textit{Student Models}}}\\
\specialrule{\lightrulewidth}{0pt}{0.8ex}
\multirow{3}{*}{InternVL2.5-4B}& Zero-Shot   & 8.8755 & 8.9423 &59.3 & 82.1  &  - \\
&LoRA & 3.4614 &3.8151& 62.3 & 87.8 &  3.7M  \\
&Full Params   & 3.0238& 3.2578& 63.4 & 90.7 & 3.0B\\
\midrule
\multirow{3}{*}{Qwen2.5-VL-3B}&Zero-Shot   & 8.0909 &8.1742& 61.2 & 85.8 &  - \\
&LoRA & 3.3101 &3.6461& 64.5 & 92.2 &  3.7M  \\

&Full Params   & 2.9127& 3.1017& 64.7 & 93.2 & 3.0B \\
\midrule
\multirow{3}{*}{Qwen2-VL-2B}& Zero-Shot   & 9.3210 &9.3872 & 58.7 & 80.4 &  - \\
&LoRA & 3.8127 & 4.3093& 61.3 & 86.8 & 4.4M   \\
&Full Params   &3.1846& 3.5204& 62.4 & 88.9 & 1.5B \\
\bottomrule
\end{tabular}
}
\label{tab:ab3}
% \vspace{-0.10 in}
\end{table*}


%只开qk 对于InternVL和Qwen2.5而言是3.7M，对于2VL2B而言是2.2M (r改32-4.4M)
%开qkvo，对于InternVL和Qwen2.5而言是7.4M，对于2BL2B而言是4.4M(r改32=8.7)




\subsection{Ablation Study}
\noindent \textbf{Component-level ablations.} Table~\ref{tab:ab1} validates the contributions of VideoStir's key components. Replacing the intent–relevance scorer with similarity-based frame retrieval results in substantial drops in both overall performance and retrieval accuracy, even when using strong vision–language embeddings from a powerful encoder such as PE. This indicates that pure semantic similarity matching can overlook the critical evidence required for long-video understanding, underscoring the importance of intent-aware retrieval. Meanwhile, removing the probability-weighted expectation (i.e., directly using discrete relevance levels) significantly degrades performance, suggesting that distributional scoring provides a smoother and better-calibrated ranking signal for discriminating frames. Moreover, disabling the spatio-temporal graph (i.e., flattening cues instead of performing structured retrieval), or removing spatial/temporal edges, markedly reduces retrieval accuracy, confirming that multi-hop aggregation over structured contexts is essential for extending evidence beyond the semantically matched anchor clips. In addition, replacing the event-boundary detector with uniform sampling that produces the same number of clips slightly degrades performance, suggesting that event-boundary segmentation isn’t the main performance driver, but it still improves evidence retrieval accuracy.

\noindent \textbf{Embedding backbones.} Table~\ref{tab:ab2} examines the impact of embedding choices on graph-based clip retrieval. Stronger video--language encoders yield consistent improvements on downstream performance. However, these gains are comparatively modest relative to the contributions from the framework’s core components, suggesting that VideoStir’s effectiveness is driven primarily by its structured design rather than by reliance on increasingly powerful embedding models.

\noindent \textbf{Scorer training strategy.} Table~\ref{tab:ab3} compares zero-shot inference, LoRA tuning, and full-parameter fine-tuning for the intent-relevance scorer. We evaluate different student-teacher pairs, including InternVL2.5~\cite{chen2024far}, Qwen2.5-VL~\cite{bai2025qwen2}, and Qwen2-VL~\cite{wang2024qwen2}. For Qwen2-VL-2B, we increase the LoRA rank to 32 (vs. 16 for others) to align trainable parameter counts for fair comparison. 

In the zero-shot setting, student's performance significantly behind teachers. Conversely, LoRA tuning markedly reduces validation cross-entropy on IR-600K and consistently improves LV-Bench performance, demonstrating robust generalization. Notably, LoRA captures the majority of full-param tuning gains using only about 4M parameters, in contrast to the billions required for full updates, with the downstream performance of LoRA-tuned Qwen2.5-VL-3B nearly matching the fully fine-tuned model. Consequently, we adopt the Qwen2.5-VL family for both teacher and student roles to maximize performance. These results confirm that IR-600K provides effective, transferable supervision, while LoRA offers an optimal accuracy–efficiency trade-off for practical long-video RAG systems.

% LVBench的超长视频VS超大窗口如gemini 2.5 flash进行视频理解？(看下gemini 2.5的官方视频理解调用文档，复现一下，持保留看法，写出来感觉会一定程度上弱化主旨，考虑等rebuttal如果遇到了再做。)

% Faliure的讨论，Fw1rirubXiU，受限于模型本身的视觉理解能力，对颜色的分类错误。

% good Case study：Z6Hx_BAZeUw_1,7分36s

%推理时间的比较和所需caption的LLM tokens，可以放到limitations里面，+附录。

%数据集严谨性

\subsection{Discussion}
Tables~\ref{tab:comp1} and~\ref{tab:comp2} present a comprehensive comparison between VideoStir and SOTA baselines, including those that utilize external expert modules (e.g., captioners, object detectors, and OCR models) to augment video with auxiliary text. Our results indicate that VideoStir achieves superior performance using only video frames as native inputs, effectively bypassing the need for external textual augmentation. This suggests that one of the primary bottlenecks in long-video understanding still lies in the retrieval and organization of pivotal visual evidence. Therefore, refining the selection of visual evidence is essential for scaling MLLM's long-video understanding capabilities. VideoStir addresses this by enhancing both the relevance and spatio-temporal coherence of retrieved evidence, thereby facilitating more reliable reasoning in downstream MLLMs.

\nocite{gu2023orsi,gu2025optical,gu2025acl,gu2025sfir,gu2025efficient,gu2024mixed,fu2025brainvis,fu2023sgcn,fu2024dp,gu2025hdtcnet,fu2025vistawise,fu2025contextnav,ren2025wamo,fu2025sdr,liu2025correlationcausationmaxpoolingbasedmultiinstance,liu2025hacsurv,zhang2024defending,zhang2025tuning,zhang2026test,ti2025towards,zhang2025tokenswap,zhang2025improving,mei2026gateddifferentiableworkingmemory,mei2025surveycontextengineeringlarge,mei2024hiddenguard,mei2025a1,mei2024not,li2024vulnerability,li2024think,li2024drs,li2025texture,xiong2025thinking,xiong2025unveiling,ge2025focusingcontrastiveattentionenhancing,ge2024can,ge2025innate,ge2025mrfdmultiregionfusiondecoding}

% Across benchmarks and backbones, the improvements correlate strongly with Retrieval Acc., reinforcing that long-video QA failures are often retrieval-driven: missing a small set of pivotal frames can invalidate downstream reasoning regardless of the generator’s capacity. VideoStir addresses this via a complementary “structure + intent” design: the spatio-temporal graph mitigates long-range context fragmentation by propagating from anchors to supported neighbors, while the intent-relevance scorer suppresses semantically similar but causally irrelevant frames and preserves implicit intent-aligned cues. The pronounced gap between distributional scoring and discrete-level selection further indicates that fine-grained calibration (rather than hard labels) matters when ranking dense, highly correlated frame streams. Finally, the LoRA-vs-SFT study suggests a practical recipe: keep the generator fixed, and invest minimal adaptation budget into a dedicated intent-aware scorer to obtain most of the retrieval and end-task gains.
%这里可以讨论一下成本、时延等优势，比如无需LLM推理成本，时延更少之类的。

\section{Conclusion}
We propose VideoStir, a novel long-video RAG framework designed to address the limitations of contextual structure decoupling and intent misalignment prevalent in current paradigms. VideoStir models long videos as spatio-temporal graphs and utilizes multi-hop traversal to recover the intrinsic topology of their contexts, thereby aggregating spatio-temporally coherent evidence. Additionally, we introduce an intent-relevance scorer trained on our newly curated IR-600K dataset, which further filters visual content by capturing the query's underlying reasoning intent beyond surface-level semantics. Both the dataset and the scorer provide a solid foundation for future research into intent-oriented long-video RAG. Extensive experiments show that VideoStir is competitive with SOTA baselines without relying on auxiliary information, underscoring the effectiveness of shifting from flattened semantic matching to structured, intent-aware reasoning, and further revealing that one of the main bottlenecks in long-video understanding lies in retrieving and organizing pivotal visual evidence.

